\documentclass[10pt,journal,compsoc,onecolumn]{IEEEtran}
\usepackage{fontspec}
\usepackage{amsmath,amsfonts,amssymb}
\usepackage{listings}
\usepackage{xcolor}

\setmainfont{Noto Serif Telugu}[
  Script=Telugu,
  BoldFont=Noto Serif Telugu Bold,
  ItalicFont=Noto Serif Telugu,
  BoldItalicFont=Noto Serif Telugu Bold
]
\newfontfamily\englishfont{DejaVu Serif}
\newfontfamily\codefont{DejaVu Sans Mono}

\lstset{
  basicstyle=\small\ttfamily,
  keywordstyle=\color{blue}\bfseries,
  frame=single
}

\begin{document}
\title{జావా టైప్ సిస్టమ్ మరియు జెనెరిక్స్ మాస్టర్‌క్లాస్: సమగ్ర విశ్లేషణ}
\maketitle

\section{పరిచయం (Introduction)}
జావా లాంగ్వేజ్ స్పెసిఫికేషన్ (JLS) ప్రకారం డేటా టైపులను రెండు విభాగాలుగా విభజించవచ్చు: 
\textbf{Primitive Types} మరియు \textbf{Reference Types}.

గణిత సమీకరణం:
\begin{equation}
S \le T \implies S[] \le T[] \quad \text{కానీ} \quad S \le T \nRightarrow \text{List}\langle S \rangle \le \text{List}\langle T \rangle
\end{equation}

\begin{lstlisting}[language=Java]
public class Box<T> {
    private T value;
    public Box(T value) { this.value = value; }
    public T getValue() { return value; }
}
\end{lstlisting}

\end{document}
